\section{Evolutionary Code and Algorithm Discovery}
\label{sec:evolutionary-code}

The preceding sections examined how individual AI components---prompts, reasoning chains, and training procedures---can improve themselves through iterative feedback. We now turn to a complementary paradigm: the use of evolutionary computation, augmented by large language models, to discover and optimize algorithms, code, and mathematical constructions at scale. This section traces the arc from treating LLMs as optimizers that operate on natural language, through LLM-driven evolutionary operators that mutate and recombine programs, to full-scale algorithm discovery systems that have produced results surpassing decades of human mathematical effort.

The central insight uniting these systems is that LLMs bring \emph{semantic understanding} to evolutionary processes that traditionally operated blindly at the syntactic level. Where classical genetic algorithms mutate bitstrings or parse trees without understanding what those modifications mean, LLM-enhanced systems propose variations that are informed by internet-scale knowledge of programming patterns, mathematical relationships, and problem structure. This semantic advantage enables a qualitative leap in the efficiency and scope of evolutionary search, transforming it from a brute-force exploration mechanism into an intelligent discovery process.

We organize this section around five major themes: (1) LLMs as general-purpose optimizers operating on natural language (OPRO), (2) LLMs as replacements for traditional evolutionary operators (OpenELM, ReEvo, EvoPrompt), (3) full-scale algorithm discovery through program evolution (FunSearch), (4) industrial-grade codebase evolution (AlphaEvolve), and (5) the open-source ecosystem that democratizes these capabilities. We conclude with a path analysis tracing the evolution from simple prompt optimization to genuine algorithm discovery.

\subsection{LLMs as Optimizers: The OPRO Paradigm}
\label{sec:opro}

Traditional optimization relies on gradient-based methods or hand-crafted heuristics, both of which require mathematical formulations of the objective and access to the internal structure of the function being optimized. OPRO~\citep{yang2023opro} (Optimization by PROmpting), introduced by Google DeepMind in 2023, takes a fundamentally different approach: it treats the large language model itself as the optimizer. Rather than computing gradients through a differentiable objective, OPRO frames optimization as a natural language reasoning task in which the LLM reads a history of candidate solutions paired with their scores, then proposes new candidates expected to achieve higher scores.

\subsubsection{Core Mechanism}

The OPRO optimization loop proceeds as follows. Given an optimization problem, the system maintains a sorted archive $\mathcal{H} = \{(s_1, f(s_1)), \ldots, (s_n, f(s_n))\}$ of instruction-score pairs, ordered from lowest to highest score. At each optimization step, the optimizer LLM receives a meta-prompt containing the top-$k$ scoring instructions from this archive, along with a set of few-shot task examples. The LLM is instructed to generate a new instruction that differs from all previous candidates and achieves a higher score. Formally:

\begin{equation}
s_{n+1} = \text{LLM}_{\text{opt}}(\mathcal{P}_{\text{meta}}(\mathcal{H}_{\text{top-}k}, \mathcal{D}_{\text{few-shot}}))
\end{equation}

where $\mathcal{P}_{\text{meta}}$ is the meta-prompt template and $\mathcal{D}_{\text{few-shot}}$ provides task context. The new instruction is then evaluated by a separate scorer LLM (with temperature set to zero for deterministic scoring) on the training set:

\begin{equation}
f(s_{n+1}) = \frac{1}{|\mathcal{T}_{\text{train}}|} \sum_{(q,a) \in \mathcal{T}_{\text{train}}} \mathbb{1}[\text{LLM}_{\text{scorer}}(s_{n+1}, q) = a]
\end{equation}

The result is inserted into the sorted archive, maintaining a maximum of 20 entries (the oldest low-scoring entries are evicted when capacity is reached).

A critical design choice is \textbf{score discretization}: continuous accuracy values are mapped into 100 discrete bins via the \texttt{\_bucketize\_float} function, preventing the optimizer from overfitting to noise in small score differences. This discretization acts as a form of quantization that stabilizes the optimization landscape, making it easier for the LLM to identify meaningful score improvements over statistical noise.

The default hyperparameters configure 200 search steps, 8 instructions generated per step, an optimizer temperature of 1.0 (high, to encourage diverse proposals), a scorer temperature of 0.0 (deterministic evaluation), and a maximum of 20 instructions retained in the meta-prompt. The high optimizer temperature is essential: it provides the stochastic exploration that replaces the random perturbations of traditional optimization algorithms, while the archive mechanism provides the exploitation pressure by always presenting the best-known solutions.

\subsubsection{Two-Layer Prompt Architecture}

OPRO employs a carefully designed two-layer prompt architecture that separates the optimization context from the task execution context.

\textbf{Outer layer (meta-prompt for the optimizer LLM)}: This prompt presents the optimization history in a structured format. Each previous instruction is shown alongside its numerical score, sorted from lowest to highest so that the LLM reads the best solutions last (a recency bias exploitation). The meta-prompt includes the instruction:
\begin{quote}
\textit{Generate an instruction that is different from all the instructions above, and has a higher score.}
\end{quote}
This explicit instruction to differentiate from previous candidates encourages exploration rather than mere reproduction of existing solutions.

\textbf{Inner layer (task prompt for the scorer LLM)}: The candidate instruction is wrapped around the actual problem instance. OPRO evaluates four possible positions for the instruction: before the question (\texttt{before\_Q}), at the question start (\texttt{Q\_begin}), at the question end (\texttt{Q\_end}), and at the answer start (\texttt{A\_begin}). The optimal position is determined empirically per dataset, with \texttt{A\_begin} typically performing best for mathematical reasoning tasks.

\textbf{Error-driven few-shot selection}: Rather than using random few-shot examples, OPRO implements a sophisticated curriculum strategy. The system tracks \texttt{wrong\_questions\_from\_start\_counter}, which accumulates the frequency of errors across all optimization steps. Few-shot examples are selected from the most frequently missed questions (\texttt{accumulative\_most\_frequent} strategy), creating a curriculum that focuses the optimizer on its weakest cases. This error-driven selection is analogous to hard negative mining in contrastive learning, directing computational resources toward the most informative training signal.

\subsubsection{Cross-Domain Generality}

A key contribution of OPRO is demonstrating that the same optimization framework applies across qualitatively different domains, requiring only changes to the natural language problem description:

\textbf{Prompt optimization}: On GSM8K (grade school mathematics), OPRO-discovered prompts achieve 80\% accuracy, up from 62\% with human-written prompts. The most famous discovered prompt---``Take a deep breath and work on this problem step-by-step''---has been widely cited as evidence that LLM-optimized instructions can outperform human intuition. On BIG-Bench Hard (27 challenging subtasks), OPRO achieves consistent improvements across diverse reasoning patterns including logical deduction, navigation, and symbolic reasoning.

\textbf{Continuous optimization}: For linear regression $f(w,b) = \|y - (Xw + b)\|^2$, the meta-prompt presents $(w, b)$ pairs with their objective values, and the LLM proposes new parameter vectors. Despite lacking gradient information, the LLM converges to near-optimal solutions within 500 steps, demonstrating that language models can perform numerical reasoning from textual score histories.

\textbf{Combinatorial optimization}: For the Traveling Salesman Problem (TSP), the LLM generates candidate routes in trace format (\texttt{<trace>0,1,2,3,\ldots</trace>}), with routes evaluated by total distance. The system discovers solutions competitive with nearest-neighbor and farthest-insertion initialization heuristics, showing that the optimization paradigm extends to discrete combinatorial problems.

\subsubsection{Convergence Analysis}

OPRO does not provide formal convergence guarantees: there is no proof that the optimization loop will find global optima or even monotonically improve. The optimization quality is bounded by the LLM's ability to reason about numerical patterns from textual score histories---a task for which language models are not inherently designed. Empirically, the system shows monotonic improvement in the majority of runs, but the trajectory is noisy and dependent on the quality of the base LLM. The score discretization mechanism partially mitigates this noise, but cannot eliminate it entirely.

The search is limited to the instruction space and cannot modify model weights, architecture, or the evaluation procedure itself. Each optimization step requires multiple LLM calls (one generation call plus $N$ scoring calls, where $N$ is the training set size), making the process expensive. On GSM8K with 7,473 training examples, a single optimization step requires roughly 7,481 LLM calls (8 generation + 7,473 scoring).

\subsubsection{Architectural Implementation}

The OPRO codebase is organized around a clean separation of concerns. The core evolutionary loop resides in \texttt{opt\_utils.py} (\texttt{run\_evolution} function, approximately 700 lines), which implements the main optimization loop, meta-prompt generation, instruction extraction, deduplication (MD5 hashing), and score computation. The evaluation framework in \texttt{eval\_utils.py} handles instruction evaluation with configurable positions, while \texttt{metrics.py} implements answer normalization supporting numerical types (unit stripping, English number conversion), boolean types (yes/no equivalence), and text types (punctuation-insensitive matching). The LLM interface layer (\texttt{prompt\_utils.py}) provides unified access to both OpenAI and Google PaLM APIs, with GPT models using \texttt{<INS>} tags and PaLM models using bracket delimiters for instruction extraction.

A GitNexus analysis reveals 697 symbol nodes, 826 dependency edges, and 9 module communities, confirming a well-modularized architecture with clear separation between optimization logic, evaluation, and LLM interfacing.

\subsubsection{Experimental Results}

Table~\ref{tab:opro-results} summarizes OPRO's performance across the evaluated benchmarks. The results demonstrate consistent improvements over baselines, though the magnitude varies significantly across task types.

\begin{table}[t]
\centering
\caption{OPRO optimization results across benchmark tasks. ``Baseline'' uses human-written prompts; ``OPRO'' uses the best prompt discovered by the optimization loop. All results use PaLM 2-L as both optimizer and scorer unless otherwise noted.}
\label{tab:opro-results}
\small
\begin{tabular}{@{}lcccc@{}}
\toprule
\textbf{Benchmark} & \textbf{Task Type} & \textbf{Baseline} & \textbf{OPRO} & \textbf{$\Delta$} \\
\midrule
GSM8K & Math reasoning & 62\% & \textbf{80\%} & +18pp \\
BBH (avg) & Reasoning (27 tasks) & Varied & \textbf{+2--12pp} & Variable \\
MMLU (avg) & Knowledge (57 tasks) & Varied & \textbf{+1--5pp} & Variable \\
Linear Regression & Continuous optimization & Suboptimal & \textbf{Near-optimal} & Converges \\
TSP (small) & Combinatorial & NN heuristic & \textbf{Competitive} & Comparable \\
\bottomrule
\end{tabular}
\end{table}

The most notable result remains the ``Take a deep breath and work on this problem step-by-step'' prompt discovered on GSM8K, which has been widely adopted beyond the research community. Analysis of the optimization trajectories reveals that the best prompts often incorporate meta-cognitive instructions (``think carefully,'' ``verify your answer'') and decomposition strategies (``break the problem into steps,'' ``work incrementally''), suggesting that the LLM optimizer is learning to construct prompts that elicit systematic reasoning from the scorer model.

On BBH, performance varies across the 27 subtasks. Tasks that benefit most from OPRO optimization tend to be those where the default prompt does not clearly specify the reasoning strategy (e.g., logical deduction, object tracking), while tasks with inherently ambiguous evaluation criteria show smaller improvements. This variation highlights a limitation of the prompt-only optimization approach: when the bottleneck is the model's fundamental capability rather than the prompt's clarity, optimization yields diminishing returns.

\subsection{LLMs as Evolutionary Operators}
\label{sec:llm-evolutionary}

While OPRO uses the LLM as a single-step optimizer within a sequential framework, a parallel line of work treats LLMs as replacements for traditional evolutionary operators---mutation, crossover, and selection---within population-based search frameworks. The key insight is that LLMs bring semantic understanding to operators that traditionally operated blindly at the syntactic level, enabling meaningful program modifications that respect functional constraints.

\subsubsection{The OpenELM Framework}

OpenELM~\citep{openelm2023} (Open-source Evolution through Large Models), developed by CarperAI (a Stability AI subsidiary), provides a general-purpose framework for LLM-driven evolutionary search. Released under the MIT license, it makes LLM-enhanced evolution accessible to the broader research community.

\textbf{Architecture}: OpenELM follows a four-layer hierarchy:
\begin{enumerate}
\item \textbf{ELM} (main controller, \texttt{elm.py}): Orchestrates the overall evolution process, managing the population and coordinating the evolution loop.
\item \textbf{MAP-Elites} (selection algorithm, \texttt{mapelites.py}): Maintains a quality-diversity grid that maps solutions to behavioral niches, ensuring that the population spans a diverse set of high-quality solutions rather than converging to a single optimum.
\item \textbf{Environment} (fitness evaluation): Problem-specific evaluation functions that compute both fitness (quality) and behavioral descriptors (for niche assignment). OpenELM provides five built-in environments: Sodarace (2D physics simulation), image evolution, programming puzzles, prompt evolution, and poetry generation (QDAIF---Quality Diversity with AI Feedback).
\item \textbf{MutationModel} (LLM variation, \texttt{mutation\_model.py}): Generates candidate modifications through LLM prompting, implementing both mutation and crossover as language-driven operations.
\end{enumerate}

\textbf{LLM-driven evolutionary operators}: For mutation, the LLM receives the parent program with the instruction: ``Modify this solution while preserving functionality but trying a different approach.'' For crossover, the LLM receives two parent programs and is asked: ``Combine the strengths of both solutions into a new solution.'' These semantic operators are qualitatively different from traditional genetic crossover, which typically operates on token or AST-node boundaries without understanding program meaning. The LLM can identify and preserve functional units, merge algorithms from different parents, and introduce novel constructs---capabilities that are beyond syntactic recombination.

\textbf{Quality-diversity algorithms}: OpenELM supports three variants of quality-diversity search:
\begin{itemize}
\item \textbf{MAP-Elites}: Standard grid-based mapping of solutions to behavioral niches, where each niche retains only the highest-performing solution.
\item \textbf{CVT-MAP-Elites}: Continuous Voronoi tessellation replaces the discrete grid, enabling smoother coverage of the behavioral space.
\item \textbf{Deep Grid MAP-Elites}: Uses deep learning to learn behavioral descriptors, enabling evolution in high-dimensional behavioral spaces where manual descriptor design is impractical.
\end{itemize}

\textbf{Safety and infrastructure}: All generated code executes within gVisor sandbox containers, providing kernel-level isolation that prevents malicious or buggy code from affecting the host system. The LangChain integration supports virtually all LLM backends (OpenAI, HuggingFace, local models), and Hydra configuration enables all parameters to be overridden from the command line.

\textbf{Built-in environments}: The five provided environments demonstrate the breadth of OpenELM's applicability. The Sodarace environment evolves 2D physics-based robots that locomote across terrain, with fitness measured by distance traveled and diversity measured by morphological features. The image evolution environment evolves visual patterns toward aesthetic criteria, with fitness measured by a pre-trained quality model and diversity measured by pixel-level features. The programming puzzles environment evolves solutions to coding challenges, with fitness measured by test case pass rates and diversity measured by solution approach. The prompt evolution environment directly optimizes LLM prompts for downstream task performance, treating the prompt as the evolvable artifact. The poetry environment (QDAIF) evolves creative text with quality measured by an LLM judge and diversity measured by text features.

Each environment follows a consistent interface: an \texttt{evaluate} method that returns fitness and behavioral descriptors, and a \texttt{to\_phenotype} method that converts the genotypic representation (typically text or code) into a phenotypic representation that can be evaluated. This interface makes it straightforward to add new environments by implementing these two methods.

The MAP-Elites implementation in OpenELM is particularly noteworthy for its flexibility. The grid can be defined with arbitrary behavioral descriptors (continuous or discrete), and the CVT-MAP-Elites variant uses $k$-means clustering to automatically discover behavioral niches rather than requiring manual specification. This automation is important for domains where the relevant behavioral dimensions are not known in advance---the system can discover them through clustering of the evolved population.

A GitNexus analysis reveals a substantial codebase with 3,596 symbol nodes, 6,639 edges, and 157 clusters, reflecting the framework's comprehensive scope across environments, algorithms, and LLM integrations.

\subsubsection{ReEvo: Reflective Evolution}

ReEvo~\citep{reevo2024} (NeurIPS 2024) extends the LLM-as-evolutionary-operator paradigm with a dual-layer reflection mechanism that mimics human expert design processes. Beyond standard mutation and crossover, ReEvo introduces:

\textbf{Short-term reflection}: After evaluating a candidate heuristic, the LLM analyzes its performance, identifying specific strengths and weaknesses. This feedback directly informs the next mutation, creating a tight optimization loop where each candidate benefits from the analysis of its predecessor.

\textbf{Long-term reflection}: Across iterations, the system accumulates structured domain knowledge in a growing text buffer. This buffer serves as a persistent memory of ``what works and what doesn't,'' accelerating convergence on subsequent iterations by preserving lessons that would otherwise be lost between runs.

ReEvo frames its approach as \emph{language hyper-heuristics} (LHHs): rather than searching within a fixed set of low-level heuristics (as traditional hyper-heuristics do), LHHs leverage the LLM's internet-scale knowledge to explore an open-ended space of heuristic designs. The framework operates in both white-box settings (where problem structure is visible to the LLM) and black-box settings (where only evaluation scores are available), demonstrating flexibility across problem modalities.

On combinatorial optimization benchmarks---TSP, CVRP, OP, MKP, BPP, and DPP---ReEvo discovers heuristics that outperform human-designed baselines, with the slogan ``5 minutes for a SOTA algorithm'' capturing its efficiency. The modular design allows new problems to be integrated by defining three components: problem configuration, evaluation pipeline, and prompt templates.

\textbf{Long-term reflection as accumulated knowledge}: The long-term reflection buffer in ReEvo deserves particular attention as a mechanism for open-ended learning. Unlike standard evolutionary algorithms that only propagate information through genetic inheritance (parent to child), ReEvo's persistent knowledge buffer accumulates \emph{domain-level insights} that transcend individual candidate solutions. For example, after exploring hundreds of TSP heuristics, the buffer might contain observations such as ``nearest-neighbor initialization tends to produce shorter initial tours than random insertion for dense graphs'' or ``2-opt local search is most effective when applied to edges that differ significantly in angle from their neighbors.'' These insights, expressed in natural language, provide the LLM with distilled expertise that accelerates subsequent search by directing attention toward the most promising regions of the heuristic space.

The interplay between short-term and long-term reflection creates a two-timescale learning process: fast adaptation within each optimization run (short-term) and slow knowledge accumulation across runs (long-term). This structure is analogous to the relationship between within-lifetime learning and evolutionary adaptation in biological systems, and it suggests that LLM-driven evolution can exhibit a form of cultural transmission where knowledge is passed between optimization runs through the persistent text buffer.

\subsubsection{EvoPrompt: Evolutionary Prompt Optimization}

EvoPrompt~\citep{evoprompt2023} applies evolutionary algorithms specifically to the prompt optimization problem. Rather than using the LLM as the optimizer (as OPRO does), EvoPrompt uses traditional evolutionary operators (mutation and crossover) applied to prompt text, with an LLM serving as the evaluation function. The mutation operator modifies a single prompt (e.g., rephrasing, adding constraints, changing the format), while the crossover operator combines elements from two high-performing prompts (e.g., merging instruction styles from one with task descriptions from another).

On text classification benchmarks, EvoPrompt discovers prompts that outperform both human-written baselines and gradient-based prompt tuning methods, demonstrating that population-based search is a viable alternative to continuous prompt optimization. The approach is particularly effective for tasks where the optimal prompt structure is non-obvious and where small wording changes can produce large performance differences.

\subsubsection{LLaMEA: Evolving Optimization Algorithms}

LLaMEA~\citep{llamea2024} (Large Language Model Evolutionary Algorithm) takes the LLM-as-evolutionary-operator concept to its logical extreme: instead of evolving solutions to specific problems, it evolves entire optimization algorithms. The system maintains a population of algorithm implementations (expressed as Python functions), each evaluated on a suite of benchmark optimization problems. The LLM proposes modifications to algorithm code---changing selection pressure, introducing new variation operators, adjusting parameter update rules---and quality-diversity search ensures that the discovered algorithms span a diverse set of optimization strategies.

This work demonstrates that LLM-driven evolution can operate at a higher level of abstraction than solution-space search: it discovers the \emph{methods} of optimization, not just the \emph{results} of optimization. The discovered algorithms sometimes incorporate strategies that resemble known optimization techniques (suggesting the LLM's training data includes relevant knowledge) and sometimes introduce genuinely novel mechanisms that do not correspond to any previously published method.

\subsection{FunSearch and Mathematical Discovery}
\label{sec:funsearch}

The systems discussed so far optimize prompts or discover heuristics for known problems. A more ambitious goal is to use LLM-driven evolution to make genuinely new mathematical discoveries. FunSearch~\citep{romera2024funsearch}, published in \emph{Nature} in 2023 by Google DeepMind, achieved exactly this: the first instance of an LLM-based system making novel scientific contributions that advanced the state of mathematical knowledge.

\subsubsection{Searching Function Space, Not Solution Space}

The core architectural insight of FunSearch is to search the space of \emph{programs that generate solutions}, rather than the space of solutions directly. This distinction is critical for several reasons. First, solution spaces for combinatorial problems are typically enormous and unstructured---the space of all possible cap sets in $\mathbb{F}_3^n$ is exponentially large and lacks the gradient structure that optimization algorithms exploit. Second, program spaces admit semantic operations that the LLM can meaningfully navigate: the LLM can understand what a function \emph{does}, not just what its output \emph{is}. Third, programs are composable and modular, allowing the evolution to discover reusable abstractions.

Concretely, FunSearch operates on Python functions decorated with two markers:
\begin{itemize}
\item \texttt{@funsearch.run}: The evaluation entry point (returns a numerical score)
\item \texttt{@funsearch.evolve}: The target function to be evolved by the LLM
\end{itemize}

This decorator-based interface is elegantly minimal: users define what to evaluate and what to evolve, and FunSearch handles the rest. The separation of evaluation from evolution ensures that the selection criteria are transparent and not influenced by the LLM's generation process.

\subsubsection{Three-Tier Population Structure}

FunSearch employs a sophisticated three-tier population structure implemented in \texttt{programs\_database.py}:

\textbf{Programs Database}: The top-level container that manages the overall evolutionary process.

\textbf{Islands}: Ten independent sub-populations that evolve in parallel. The island model is a well-established technique in evolutionary computation for maintaining diversity and preventing premature convergence. FunSearch's implementation adds a distinctive feature: every 4 hours (the \texttt{reset\_period}), the bottom 50\% of islands (ranked by best score) are reset. New founding programs are drawn from the top-performing islands, implementing a form of migration that concentrates search effort on the most promising regions while periodically introducing fresh genetic material.

\textbf{Clusters}: Within each island, programs are grouped by their \emph{Signature}---the tuple of scores across all test cases. Programs with identical signatures (i.e., programs that perform equivalently on all test inputs) are clustered together. Within each cluster, shorter programs are preferred, implementing an Occam's razor principle that encourages simplicity and discourages bloat (the tendency of evolved programs to grow without bound).

\textbf{Temperature-based sampling}: Program selection within islands uses a softmax with an initial temperature of 0.1, heavily biasing selection toward high-scoring clusters while maintaining some probability of exploring lower-scoring regions. The low initial temperature ensures that the search quickly focuses on the most promising areas, while the nonzero temperature prevents complete convergence.

\subsubsection{Prompt Engineering for Evolution}

The prompt design in FunSearch is particularly elegant and represents a significant engineering contribution. Each prompt presents the LLM with a versioned evolution history:

\begin{verbatim}
def priority_v0(element, n):
  """Returns the priority with which we want
  to add `element`."""
  priority = element
  return priority

def priority_v1(element, n):
  """Improved version of `priority_v0`."""
  priority = element ** 2
  return priority

def priority_v2(element, n):
  """Improved version of `priority_v1`."""
\end{verbatim}

Six design patterns underlie this prompt structure:
\begin{enumerate}
\item \textbf{Versioned naming}: \texttt{v0}, \texttt{v1}, \texttt{v2} suffixes make the evolution chain explicit, helping the LLM understand that it is building on prior work rather than generating from scratch.
\item \textbf{Progressive docstrings}: ``Improved version of \texttt{priority\_v1}'' provides clear semantic guidance about the task---the LLM should improve upon the existing best, not generate an arbitrary function.
\item \textbf{Score-ordered presentation}: Lower-scoring versions appear first, higher-scoring versions last. This leverages the LLM's recency bias: the last thing it reads is the best current solution, which anchors its generation.
\item \textbf{Empty function body}: The latest version (\texttt{v2}) has only a signature and docstring, inviting the LLM to complete the implementation. This provides structure without constraining the solution space.
\item \textbf{Tokenizer-level renaming}: The \texttt{rename\_function\_calls()} function uses Python's tokenizer (not string replacement) to rename function references throughout the code, ensuring that recursive or cross-referential programs are handled correctly.
\item \textbf{Preserved imports and helpers}: The prompt retains complete import statements and auxiliary functions, providing the LLM with the full context needed to generate syntactically correct and semantically meaningful code.
\end{enumerate}

This prompt engineering represents a form of \emph{meta-design}: the way the evolutionary history is presented to the LLM determines the efficiency of the search, and FunSearch's designers invested significant effort in optimizing this presentation.

\subsubsection{Evaluation Framework}

The evaluation pipeline (\texttt{evaluator.py}) implements three tiers of safety filtering:

\begin{enumerate}
\item \textbf{Execution success}: Programs must run without errors in a sandboxed environment with a 30-second timeout. This catches syntax errors, infinite loops, and runtime exceptions.

\item \textbf{Ancestor call detection}: The system checks whether a generated program directly calls its parent's implementation (which would effectively ``cheat'' by delegating to a known-good solution rather than implementing a genuine improvement). The \texttt{\_calls\_ancestor} function uses AST analysis to detect such delegations.

\item \textbf{Return type validation}: Programs must return numerical values (int or float). Non-numerical returns (strings, lists, None) are rejected.
\end{enumerate}

Code post-processing uses Python AST parsing (\texttt{\_trim\_function\_body()}) to iteratively trim malformed LLM output until it becomes syntactically valid. This robust handling of incomplete or garbled code generation is essential for maintaining a productive evolution loop---without it, many LLM outputs would be immediately rejected, wasting the generation budget.

The default configuration uses 10 islands, 2 functions per prompt, 15 samplers, 140 evaluators (a 1:9 sampler-to-evaluator ratio reflecting evaluation as the computational bottleneck), 4 samples per prompt, and a 4-hour island reset period. The 140 evaluators run in parallel, continuously processing candidate programs generated by the 15 samplers.

\subsubsection{Mathematical Discoveries}

FunSearch produced three notable mathematical results, each demonstrating the system's ability to make genuinely new contributions:

\textbf{Cap set problem}: For the problem of constructing large cap sets in $\mathbb{F}_3^n$ (sets with no three-term arithmetic progressions), FunSearch discovered new constructions exceeding the best previously known bounds. This was the first new result on this problem in over two decades, advancing a question that had resisted improvement despite significant mathematical attention.

\textbf{Bin packing}: For the online bin packing problem (a classic NP-hard optimization problem), FunSearch discovered a new heuristic that outperforms all previously known constructive heuristics on standard benchmark instances. The discovered heuristic incorporates sorting criteria that had not been previously considered in the bin packing literature.

\textbf{Cyclic graph independent sets}: FunSearch found new constructions for independent sets in cyclic graphs, contributing to the combinatorics literature on graph theory.

These results are significant not merely for their mathematical content, but for what they demonstrate about the potential of LLM-driven evolutionary search as a tool for scientific discovery. The system did not merely reproduce known results (though it did so for 75\% of test problems); it generated genuinely novel constructions that extended the frontier of mathematical knowledge.

\subsubsection{The Significance of Function-Space Search}

The choice to search function space rather than solution space has deeper implications than may be immediately apparent. When searching solution space directly, the evolutionary process must discover good solutions \emph{de novo} in each run, with no mechanism for generalizing across problem instances. When searching function space, the evolutionary process discovers \emph{methods} that can be applied to arbitrary instances of the problem, producing solutions that are inherently generalizable.

This distinction is crucial for understanding why FunSearch was able to make mathematical discoveries that prior LLM-based approaches could not. Earlier systems that used LLMs to directly solve mathematical problems (rather than evolving programs that solve them) were limited by the LLM's ability to produce correct answers for specific instances. FunSearch, by evolving the \emph{function} that generates solutions, could discover mathematical structures that no single LLM call would produce.

The function-space approach also has important implications for the scalability of evolutionary discovery. As problem sizes increase, the solution space grows exponentially, making direct search intractable. Function space, while also large, grows more slowly because functions can be expressed concisely (a short Python function can generate exponentially many solutions). This compression means that the evolutionary process operates in a lower-dimensional space than direct solution search, making it more efficient and more likely to discover generalizable patterns.

\subsubsection{Infrastructure Requirements}

FunSearch's design reveals an important architectural pattern: the separation of the evolutionary algorithm from the infrastructure required to run it. The open-source release provides the algorithmic framework (population management, island model, prompt construction, code evaluation) as a clean Python codebase, but leaves the LLM interface and sandbox execution as abstract interfaces that users must implement. This design reflects the reality that the LLM API and sandbox infrastructure are deployment-specific and evolve rapidly, while the evolutionary algorithm is more stable.

The default configuration (15 samplers, 140 evaluators) suggests that evaluation is the primary computational bottleneck, consuming roughly 90\% of the resources. This 1:9 ratio between generation and evaluation is a consequence of the asymmetry between LLM inference (relatively fast, especially with batched requests) and program execution (which may involve expensive computation, especially for mathematical problems with large test cases). Future systems may address this bottleneck through surrogate evaluation models that approximate expensive evaluations, allowing the evolutionary process to screen candidates cheaply before committing to full evaluation.

\subsection{AlphaEvolve: Full-Codebase Evolution at Scale}
\label{sec:alphaevolve}

AlphaEvolve~\citep{novikov2025alphaevolve}, also from Google DeepMind, extends the FunSearch paradigm from evolving individual functions to evolving entire codebases. Published in 2025, it represents the current state of the art in LLM-driven evolutionary code optimization, combining Gemini Flash and Gemini Pro models with MAP-Elites quality-diversity search at Google scale.

\subsubsection{Dual-Model Architecture}

The system employs two complementary LLM tiers in a breadth-depth strategy:

\textbf{Gemini Flash} (breadth): Generates many diverse candidate modifications quickly, exploring a wide range of possible improvements. The high throughput of Flash enables rapid exploration of the modification space, proposing hundreds of variations per evolution cycle.

\textbf{Gemini Pro} (depth): Analyzes promising candidates in depth, proposing sophisticated refinements that leverage deeper reasoning capabilities. Pro is reserved for candidates that show initial promise, concentrating expensive model calls where they are most likely to yield improvements.

This dual-model approach addresses a fundamental tradeoff in LLM-driven evolution: weak models propose shallow mutations at low cost, while strong models propose deep modifications at high cost. By combining both tiers, AlphaEvolve achieves the benefits of broad exploration (many diverse candidates) and deep exploitation (sophisticated improvements on promising directions).

\subsubsection{MAP-Elites Quality-Diversity Search}

The evolutionary database maintains a population of candidate programs, each annotated with fitness scores and behavioral descriptors. The MAP-Elites update rule replaces the existing elite in a behavioral niche if the child achieves a higher fitness:

\begin{equation}
\text{Archive}[b(p_{\text{child}})] = \begin{cases} p_{\text{child}} & \text{if } f(p_{\text{child}}) > f(\text{Archive}[b(p_{\text{child}})]) \\ \text{Archive}[b(p_{\text{child}})] & \text{otherwise} \end{cases}
\end{equation}

where $b(\cdot)$ maps a program to its behavioral niche descriptor. The behavioral descriptors are problem-specific features that characterize how a solution behaves (not just its quality), enabling the archive to maintain a diverse collection of high-performing solutions that approach the problem in qualitatively different ways.

Parent selection is weighted by both fitness and novelty, ensuring that underexplored regions of the behavioral space receive attention alongside high-performing regions. This dual pressure---quality and diversity---is what distinguishes MAP-Elites from simple hill-climbing or tournament selection.

\subsubsection{Results Across Domains}

AlphaEvolve's results span three categories, demonstrating broad applicability beyond mathematical discovery:

\textbf{Mathematical discoveries}:
\begin{itemize}
\item Found a $4\times 4$ complex matrix multiplication algorithm using only \textbf{48 multiplications}, the first improvement over Strassen's algorithm in 56 years. Strassen's 1969 algorithm achieved 49 multiplications for this case, and despite decades of effort by mathematicians, no improvement had been found.
\item Re-discovered state-of-the-art solutions for 75\% of 50+ mathematical problems from diverse areas, validating the system's generality.
\item Found improvements over existing best-known results for 20\% of the test problems, generating genuinely new knowledge.
\end{itemize}

\textbf{Google infrastructure improvements}:
\begin{itemize}
\item Data center scheduling (Borg system): Recovered \textbf{0.7\%} of worldwide Google compute capacity. Given Google's massive infrastructure, this seemingly small percentage translates to enormous resource savings.
\item Hardware accelerator design: Achieved functionally equivalent designs with simplified implementations, reducing design complexity.
\end{itemize}

\textbf{LLM training optimization}:
\begin{itemize}
\item FlashAttention kernel tiling: Achieved a \textbf{23\% speedup} in attention computation, directly improving the efficiency of LLM training and inference.
\item Attention computation: Achieved a \textbf{32\% speedup} in another attention kernel, demonstrating that LLM-driven evolution can optimize the very systems that enable it.
\end{itemize}

These results demonstrate that LLM-driven evolutionary search can produce practical value at industrial scale. The infrastructure improvements are not academic curiosities---they represent real efficiency gains deployed across Google's production systems.

\subsubsection{Comparison with FunSearch}

AlphaEvolve differs from FunSearch in several important ways:

\textbf{Scope}: AlphaEvolve operates on entire codebases rather than individual functions, expanding the search space from single-function optimization to multi-file, multi-module program evolution. This expanded scope enables the discovery of optimizations that span module boundaries---the kind of cross-cutting improvements that are invisible to function-level evolution. For example, an AlphaEvolve optimization might co-optimize the data structure layout in one module with the access pattern in another, producing a combined improvement that neither optimization could achieve alone.

\textbf{Model architecture}: The dual-model approach (Flash + Pro) provides a breadth-depth tradeoff that FunSearch's single-model setup cannot replicate. FunSearch uses a single LLM for all generation, while AlphaEvolve allocates model capacity based on the potential impact of each modification. The Flash model's high throughput enables rapid exploration of the modification space, while the Pro model's deeper reasoning enables sophisticated refinements on the most promising candidates.

\textbf{Target domains}: AlphaEvolve targets real-world infrastructure optimization alongside mathematical discovery, demonstrating broader applicability. FunSearch focused exclusively on mathematical problems with clean evaluation criteria, while AlphaEvolve extends to production systems where the evaluation criteria may be more complex (involving multiple objectives such as speed, memory usage, and correctness).

\textbf{Accessibility}: Both systems share a fundamental dependency on well-defined evaluation functions. Where mathematical problems have clear correctness criteria, infrastructure optimization requires carefully designed objective functions that may not capture all relevant considerations (maintainability, readability, edge cases not covered by the evaluation suite). The design of these objective functions is itself a form of engineering that requires domain expertise and careful consideration of what the optimization should prioritize.

\subsubsection{The 0.7\% Compute Recovery in Context}

The Borg data center scheduling improvement deserves special analysis. Recovering 0.7\% of worldwide Google compute capacity may sound modest, but given Google's infrastructure scale---millions of servers across dozens of data centers---this translates to resources equivalent to hundreds of thousands of additional machines. The improvement was achieved by evolving the scheduling algorithm that assigns tasks to machines, optimizing for resource utilization while respecting constraints on latency, availability, and power consumption.

What makes this result particularly notable is that the scheduling algorithm had already been heavily optimized by human engineers over many years. The fact that AlphaEvolve could find further improvements suggests that LLM-driven evolutionary search can explore regions of the algorithm design space that human engineers do not naturally consider, either because the modifications are too complex to reason about manually or because they violate engineering intuitions that happen to be suboptimal.

The FlashAttention kernel optimization (23\% speedup) and attention computation optimization (32\% speedup) are similarly significant because they directly improve the efficiency of LLM training and inference---the very systems that enable AlphaEvolve itself. This self-referential quality, where the evolved code improves the infrastructure that enables the evolution, points toward a potential virtuous cycle: more efficient LLM training enables more powerful evolutionary search, which discovers further optimizations, and so on.

\subsubsection{Limitations}

AlphaEvolve is not publicly available (academic early access only), requiring access to Google's Gemini models and compute infrastructure. The dual-model approach requires access to multiple LLM tiers, limiting reproducibility. Each evolution cycle is computationally expensive, and the evaluation functions must be carefully designed to avoid optimizing for proxy metrics that do not reflect real-world quality. The system is limited to domains with automated verifiers---creative tasks, open-ended reasoning, and domains without clear correctness criteria remain beyond its scope.

\subsection{Open-Source Ecosystem}
\label{sec:opensource-evo}

The impressive results of FunSearch and AlphaEvolve have catalyzed a growing open-source ecosystem that seeks to democratize LLM-driven evolutionary code optimization. While no open-source system can replicate Google's proprietary infrastructure (Gemini models, massive compute, internal evaluation pipelines), several projects provide accessible implementations of the core concepts.

\subsubsection{OpenEvolve}

OpenEvolve~\citep{openevolve2025} is a community-driven open-source implementation of AlphaEvolve's core concepts, available at \texttt{github.com/codelion/openevolve}. It provides the fundamental building blocks of LLM-driven codebase evolution:

\textbf{Evolutionary loop}: The main evolution cycle---parent selection, LLM-driven mutation, evaluation, and archive update---is implemented as a configurable pipeline. Users specify the evaluation function, the LLM backend, and the search parameters.

\textbf{MAP-Elites archive}: A quality-diversity archive maintains a map of high-performing solutions across behavioral dimensions, implementing the same niche-based selection pressure as AlphaEvolve.

\textbf{LLM-driven mutation}: Mutation operators are implemented as LLM prompts that present the parent program with evolution context (prior attempts, scores, task description). The system supports multiple LLM backends through a unified interface.

\textbf{Alternative backends}: Unlike AlphaEvolve (which requires Gemini models), OpenEvolve supports any LLM accessible via API, including open-source models (Llama, Mistral, Qwen) and commercial APIs (OpenAI, Anthropic). This flexibility enables researchers to experiment with different model capabilities and cost-performance tradeoffs.

OpenEvolve enables researchers to experiment with the AlphaEvolve paradigm on smaller-scale problems and with alternative LLM backends, making the approach accessible to academic labs and individual researchers.

\subsubsection{SE-Agent: Multi-Trajectory RRR}

SE-Agent implements a Multi-trajectory Revision/Recombination/Refinement (RRR) approach to code evolution. Rather than evolving a single program through sequential mutations, SE-Agent collects multiple execution trajectories---both successful and failed---and uses them collectively to guide improvement.

\textbf{Revision}: Failed trajectories are analyzed to extract error patterns. The LLM identifies common failure modes (off-by-one errors, incorrect boundary conditions, wrong algorithmic choices) and generates targeted fixes. This phase transforms failure into learning signal.

\textbf{Recombination}: Successful sub-trajectories---portions of execution traces that produced correct intermediate results---are identified and combined. The LLM merges the best elements from different successful attempts, creating novel combinations that may outperform any individual trajectory.

\textbf{Refinement}: The combined solution is iteratively improved based on aggregated evaluation feedback. Multiple evaluation criteria (correctness, efficiency, code quality) guide the refinement process.

The RRR approach is complementary to FunSearch's island model: where FunSearch maintains population diversity through independent islands, SE-Agent maintains diversity through multi-trajectory analysis. Empirically, the combination of revision and recombination produces larger improvements than either alone, suggesting that learning from both failures and successes is more effective than learning from successes alone.

\subsubsection{CodeEvolve}

CodeEvolve provides a test-driven code evolution framework where generated code is evaluated against automated test suites. The system follows an iterative cycle: generate code, execute tests, extract failure information, and use the failure feedback to guide the next generation. This approach is particularly well-suited for domains where test suites provide comprehensive coverage and where correctness can be unambiguously determined.

\subsubsection{EvoAgentX}

EvoAgentX~\citep{evoagentx2024} extends evolutionary optimization from code to multi-agent systems. Rather than evolving individual programs, EvoAgentX evolves agent configurations---their prompts, tool selections, communication patterns, and decision-making strategies. The framework combines agent-level evolution with population management, enabling the discovery of novel multi-agent coordination strategies.

\subsubsection{Comparative Analysis}

Table~\ref{tab:opensource-comparison} provides a systematic comparison of the open-source evolutionary code systems.

\begin{table*}[t]
\centering
\caption{Comparison of open-source LLM-driven evolutionary code systems. ``Search space'' indicates what the system evolves; ``Diversity'' shows the mechanism for maintaining population diversity; ``LLM backend'' lists supported language model interfaces; ``Scale'' indicates the typical problem size the system targets.}
\label{tab:opensource-comparison}
\small
\begin{tabular}{@{}llllll@{}}
\toprule
\textbf{System} & \textbf{Search Space} & \textbf{Diversity} & \textbf{LLM Backend} & \textbf{License} & \textbf{Scale} \\
\midrule
OpenEvolve & Full codebases & MAP-Elites & Multi-backend & Open & Medium \\
FunSearch & Math functions & Island model & Single (closed) & Apache 2.0 & Large \\
SE-Agent & Programs & Multi-trajectory & Multi-backend & Open & Medium \\
CodeEvolve & Programs & Elite preservation & Multi-backend & Open & Small \\
EvoAgentX & Agent configs & Population mgmt & Multi-backend & Open & Medium \\
OpenELM & Multi-domain & MAP-Elites grid & LangChain (any) & MIT & Medium \\
ReEvo & Heuristic code & Dual reflection & Multi-backend & Open & Small--Med \\
\bottomrule
\end{tabular}
\end{table*}

Several patterns emerge from the ecosystem comparison. All open-source systems use multi-backend LLM support, reflecting the community's need for flexibility in model selection. MAP-Elites and island models are the dominant diversity mechanisms, with recent systems exploring multi-trajectory and reflection-based alternatives. The ecosystem is converging on a common architectural pattern: evolution loop + quality-diversity archive + LLM mutation + automated evaluation.

\subsubsection{The Accessibility Gap}

A significant gap exists between the proprietary systems (FunSearch, AlphaEvolve) that have produced the most impressive results and the open-source systems that aim to democratize these capabilities. This gap manifests along several dimensions:

\textbf{Model access}: AlphaEvolve uses Gemini Flash and Pro, which are not publicly available at the same scale. Open-source systems must rely on publicly available models (GPT-4, Claude, or open-weight models like Llama and Mistral), which may be less capable at the specific task of proposing semantically meaningful code modifications.

\textbf{Compute scale}: Google's infrastructure enables AlphaEvolve to run thousands of evaluation episodes in parallel. Open-source systems typically run on single machines or small clusters, limiting the population sizes and evolution durations that are practically achievable.

\textbf{Evaluation infrastructure}: AlphaEvolve benefits from Google's internal benchmark suites, testing infrastructure, and domain-specific evaluation functions. Open-source systems must construct their own evaluation infrastructure, which may not be as comprehensive or well-tested.

Despite these gaps, the open-source ecosystem is rapidly maturing. OpenEvolve has demonstrated that the core AlphaEvolve paradigm can be replicated with publicly available components, albeit at smaller scale. As open-weight LLMs continue to improve and cloud compute becomes cheaper, the accessibility gap is likely to narrow. The key question is whether the most impressive discoveries (like the 48-multiplication Strassen improvement) require Google-scale resources, or whether they could be achieved with more efficient search strategies that reduce the computational requirements.

\subsection{From Code Evolution to Algorithm Discovery}
\label{sec:code-to-algorithm}

The systems surveyed in this section represent a progression of increasing ambition and capability. We can identify five levels of evolutionary complexity:

\textbf{Level 1---Prompt optimization} (OPRO): The simplest form of LLM-driven evolution optimizes natural language instructions that guide model behavior. The search space is the set of all possible natural language prompts, and the fitness function is task performance with the prompt. OPRO demonstrates that even this relatively constrained space admits significant optimization: the ``Take a deep breath'' prompt on GSM8K represents an 18 percentage point improvement over human-written baselines.

\textbf{Level 2---Code mutation} (OpenELM, ReEvo): The search space expands to include program code, with LLM-driven operators performing semantic mutations and crossovers. The key advance is that the LLM understands what the code \emph{does}, enabling modifications that preserve functionality while exploring alternative implementations. ReEvo's dual-reflection mechanism adds a learning component that accumulates domain knowledge across iterations.

\textbf{Level 3---Function evolution} (FunSearch): The search space narrows to focus on specific mathematical functions, but the ambition expands from optimization to discovery. FunSearch's function-space search (rather than solution-space search) enables the system to discover new mathematical constructions, not just optimize known ones. The island model with signature clustering provides the population diversity needed for open-ended exploration.

\textbf{Level 4---Codebase evolution} (AlphaEvolve): The search space expands dramatically to encompass entire codebases, enabling cross-module optimizations that are invisible to function-level evolution. The dual-model architecture (Flash + Pro) provides a principled approach to the breadth-depth tradeoff, and the industrial-scale results (Borg scheduling, FlashAttention optimization) demonstrate practical impact.

\textbf{Level 5---Algorithm discovery} (emerging): The frontier of the field aims to discover fundamentally new algorithms---not just optimize implementations of known algorithms, but invent algorithmic approaches that have never been conceived. AlphaEvolve's 48-multiplication Strassen improvement points toward this frontier, and future systems may discover algorithms in domains where no efficient solution is currently known.

\subsubsection{Theoretical Considerations}

The progression from Level 1 to Level 5 raises several theoretical questions about the limits of LLM-driven evolutionary search.

\textbf{Search space complexity}: The size of the search space grows dramatically across levels. Prompt spaces are effectively infinite (any natural language string is a valid prompt), but structured by linguistic constraints. Code spaces are larger still (any syntactically valid program), and codebase spaces are combinatorially explosive. Despite this growth, the quality of discovered solutions does not necessarily degrade---the LLM's semantic understanding provides search guidance that partially compensates for the larger space. This suggests that the effective dimensionality of the search (the subspace the LLM meaningfully explores) may be much smaller than the nominal dimensionality.

\textbf{Exploration-exploitation tradeoff}: All five levels face the fundamental tension between exploring new regions of the search space and exploiting known good solutions. Quality-diversity algorithms (MAP-Elites, island models) provide one resolution by explicitly maintaining diverse solutions alongside high-quality ones. The LLM's stochastic generation (via temperature sampling) provides another: high temperature encourages exploration, while the archive/history mechanism provides exploitation pressure. The optimal balance varies across levels---function evolution (Level 3) benefits from stronger exploitation (FunSearch's low temperature of 0.1), while codebase evolution (Level 4) requires more exploration to navigate the larger space.

\textbf{Grounding and verification}: A critical enabler of higher-level discovery is the availability of automated verification. Mathematical problems (Levels 3--4) have clear correctness criteria: either a matrix multiplication uses fewer scalar multiplications than Strassen's algorithm, or it does not. This binary verification provides unambiguous selection pressure. For domains without automated verification (creative writing, system design, strategy formulation), the evaluation bottleneck becomes the limiting factor, preventing progress beyond Level 2.

\textbf{The role of LLM scaling}: The quality of evolutionary discovery is correlated with the capability of the underlying LLM. AlphaEvolve's use of Gemini Pro (a frontier model) for deep analysis is not incidental---the depth of insight that the LLM can provide directly bounds the quality of proposed modifications. As LLMs continue to scale in capability, the ceiling for evolutionary discovery will rise correspondingly. However, scaling alone does not guarantee discovery: the evaluation framework, diversity maintenance, and search strategy remain critical design choices that determine whether LLM capability is effectively translated into discovery.

\subsubsection{Enablers of Algorithm Discovery}

Several factors enable the transition from optimization to discovery:

\begin{enumerate}
\item \textbf{Automated evaluation}: Mathematical verification and comprehensive test suites provide reliable selection pressure without human intervention. Without automated evaluation, the evolutionary loop cannot scale.

\item \textbf{Quality-diversity search}: MAP-Elites, island models, and other quality-diversity algorithms prevent convergence to local optima, maintaining a diverse population of high-quality solutions that can serve as stepping stones to unexpected discoveries.

\item \textbf{LLM-driven semantic variation}: Unlike random mutation, LLM-driven variation proposes modifications informed by internet-scale knowledge of programming patterns and mathematical relationships. This semantic advantage enables the search to make qualitatively different proposals than purely random exploration.

\item \textbf{Sufficient compute scale}: The most impressive results (FunSearch's cap set discovery, AlphaEvolve's Strassen improvement) require significant computational resources---hundreds of evaluators running continuously for hours or days.

\item \textbf{Well-defined search spaces}: Mathematical problems provide particularly clean search spaces with unambiguous evaluation criteria. The challenge for future work is extending algorithm discovery to domains with less clear-cut evaluation.
\end{enumerate}

\subsubsection{Challenges and Limitations}

Despite impressive progress, several fundamental challenges remain:

\textbf{Evaluation bottleneck}: All systems require well-defined evaluation functions. For mathematical problems and benchmark tasks, such functions are readily available. For open-ended creative tasks, real-world software engineering, or domains where correctness is hard to formalize, the evaluation bottleneck becomes the limiting factor. AlphaEvolve's infrastructure optimizations rely on benchmark-defined objectives that may not capture all aspects of production quality (maintainability, readability, edge case handling).

\textbf{Computational cost}: Evolutionary search with LLM-based operators is expensive. FunSearch defaults to 15 samplers and 140 evaluators running continuously; AlphaEvolve leverages Google-scale compute infrastructure. OpenELM and ReEvo reduce costs by using smaller LLMs and shorter evolution horizons, but at the cost of discovery quality. Democratizing these approaches requires more efficient search strategies---for example, surrogate models that approximate expensive evaluations, or adaptive allocation of LLM calls based on expected information gain.

\textbf{LLM capability ceiling}: The quality of evolved solutions is fundamentally bounded by the LLM's capability to propose meaningful variations. Weak LLMs produce shallow mutations; strong LLMs are expensive. The dual-model approach (AlphaEvolve's Flash+Pro) partially addresses this by allocating expensive model calls selectively, but the tradeoff remains. As LLMs continue to improve, the ceiling for evolutionary discovery will rise correspondingly.

\textbf{Safety and verification}: Self-modifying code introduces safety risks that are not fully addressed by current systems. FunSearch's three-tier filter catches obvious issues but cannot guarantee the absence of subtle bugs, security vulnerabilities, or unintended behaviors in evolved programs. As evolutionary code systems are applied to safety-critical domains (medical software, autonomous systems, financial algorithms), more rigorous verification will be essential.

\textbf{Reproducibility gap}: Both FunSearch and AlphaEvolve use proprietary infrastructure (Google's LLMs and compute) that is not publicly available. OpenEvolve and OpenELM provide open-source alternatives, but cannot replicate the scale of the original results. The gap between published results and reproducible science remains significant, and the field would benefit from standardized benchmarks and evaluation protocols that enable fair comparison across systems.

\subsubsection{Comprehensive System Comparison}

Table~\ref{tab:evo-comparison} provides a systematic comparison of the major LLM-driven evolutionary systems discussed in this section.

\begin{table*}[t]
\centering
\caption{Comparison of LLM-driven evolutionary code systems. ``Search space'' indicates what the system evolves; ``Evolution operators'' describes how variation is performed; ``Diversity mechanism'' shows how population diversity is maintained; ``LLM calls/step'' estimates the computational cost per evolution step.}
\label{tab:evo-comparison}
\small
\begin{tabular}{@{}llllcc@{}}
\toprule
\textbf{System} & \textbf{Search Space} & \textbf{Evolution Operators} & \textbf{Diversity Mechanism} & \textbf{LLM Calls/Step} & \textbf{Key Result} \\
\midrule
OPRO & Instructions & Single-step generation & Top-$k$ archive & 8 gen + $N$ eval & +18pp on GSM8K \\
OpenELM & Programs (multi-domain) & LLM mutation + crossover & MAP-Elites grid & 1--2 & Sodarace optimization \\
ReEvo & Heuristic code & Mutation + dual reflection & Elite preservation & 3--5 per candidate & SOTA on 6 combinatorial problems \\
EvoPrompt & Prompts & Mutation + crossover & Population selection & 2--4 per generation & Improved text classification \\
LLaMEA & Algorithm code & LLM mutation & Quality-diversity archive & 1--3 per candidate & Novel optimization algorithms \\
FunSearch & Math functions & LLM mutation & Island model + signature clustering & 4 per prompt, 15 samplers & New cap set bounds \\
AlphaEvolve & Full codebases & Dual-model mutation & MAP-Elites archive & Variable (Flash+Pro) & 48-mult Strassen improvement \\
OpenEvolve & Full codebases & LLM mutation & MAP-Elites & Variable & Community AlphaEvolve \\
\bottomrule
\end{tabular}
\end{table*}

Several cross-cutting patterns emerge from this comparison. Nearly all systems use some form of quality-diversity maintenance---MAP-Elites (OpenELM, AlphaEvolve, OpenEvolve), island models (FunSearch), or elite preservation (ReEvo)---reflecting a recognition that single-objective search is prone to local optima. All systems cleanly separate the LLM's role as a generator of candidates from the automated evaluation of those candidates, allowing search to scale without human evaluation. Even systems that evolve code rely heavily on carefully engineered prompts to guide the LLM's variation operators, making prompt design itself a form of ``meta-evolution'' that determines search efficiency.

\subsubsection{Convergence Toward a Common Architecture}

Despite their independent origins and different target domains, the systems surveyed in this section are converging toward a common architectural template:

\begin{enumerate}
\item \textbf{Evolution loop}: A main loop that repeatedly generates candidates, evaluates them, and updates the population or archive. This loop is the skeleton upon which all other components are hung.

\item \textbf{LLM-driven variation}: An LLM (or LLM ensemble) that proposes modifications to existing solutions. The LLM serves as an intelligent mutation operator, bringing semantic understanding to the variation process.

\item \textbf{Automated evaluation}: A programmatic evaluation function that scores candidates without human intervention. The evaluation function is the source of selection pressure, and its quality determines the ceiling for evolutionary discovery.

\item \textbf{Quality-diversity maintenance}: A mechanism (MAP-Elites, island model, archive) that maintains a diverse population of high-quality solutions, preventing convergence to local optima and enabling open-ended exploration.

\item \textbf{Prompt engineering}: Carefully designed prompts that present the LLM with evolution context (prior solutions, scores, task description, improvement history). The prompt design is itself a form of optimization that determines the efficiency of the search.
\end{enumerate}

This convergence suggests that the field is maturing toward a standardized approach, with innovations occurring primarily in the specific instantiation of each component (which LLM to use, how to design the evaluation function, how to maintain diversity) rather than in the overall architecture. The emergence of open-source frameworks (OpenELM, OpenEvolve) that provide this architecture as a reusable platform supports this interpretation.

These limitations point toward important research directions: developing more efficient evaluation methods, exploring multi-objective optimization that balances multiple quality criteria, and creating standardized benchmarks that enable fair comparison across evolutionary systems. We return to these themes in Section~\ref{sec:open-problems} when discussing future directions.

\subsubsection{Summary of Key Insights}

The evolutionary code and algorithm discovery paradigm, as surveyed in this section, yields several insights that are relevant beyond the immediate domain:

\textbf{Insight 1: Semantic variation enables discovery}. The consistent superiority of LLM-driven variation operators over random mutation demonstrates that semantic understanding of the search space is a critical enabler of evolutionary discovery. The LLM does not merely propose syntactically valid variations; it proposes variations that are \emph{meaningful} in the context of the problem, dramatically reducing the number of evaluations needed to find high-quality solutions.

\textbf{Insight 2: Quality-diversity search outperforms pure optimization}. The prevalence of MAP-Elites, island models, and other quality-diversity mechanisms across all surveyed systems reflects a fundamental insight: in open-ended search, maintaining a diverse population of high-quality solutions is more productive than optimizing a single solution. Diversity provides stepping stones to unexpected discoveries and prevents the search from becoming trapped in local optima.

\textbf{Insight 3: The evaluation function is the bottleneck}. The quality of evolutionary discovery is fundamentally bounded by the quality of the evaluation function. Mathematical problems (where evaluation is unambiguous) yield the most impressive discoveries; domains without clear evaluation criteria remain challenging. This bottleneck motivates research into more general evaluation methods, including learned evaluators and human-in-the-loop evaluation.

\textbf{Insight 4: Prompt engineering is meta-evolution}. The design of prompts that guide the LLM's variation operators is itself a form of optimization that determines the efficiency of the search. FunSearch's six prompt design patterns, OPRO's error-driven few-shot selection, and ReEvo's dual-reflection mechanism all represent engineered solutions to the meta-optimization problem of how to present evolution context to the LLM. Automating this meta-optimization---evolving the prompts that guide evolution---is a promising direction for future research.

\textbf{Insight 5: Scale matters for discovery}. The most impressive results (FunSearch's cap set discovery, AlphaEvolve's Strassen improvement, the 0.7\% Borg improvement) all require significant computational resources. This scale dependence creates a tension with the democratization goal: open-source systems cannot replicate the scale of proprietary systems, limiting the accessibility of evolutionary discovery. Surrogate evaluation models and more efficient search strategies may help bridge this gap.

The evolutionary code paradigm described in this section provides one axis of self-evolution---improving the \emph{programs} that AI systems execute. The next section examines a complementary axis: improving the \emph{agents} themselves---their architectures, workflows, and learning mechanisms.
